\hypertarget{results-of-the-bounded-final-exploration}{%
\subsection{Results of the bounded final
exploration}\label{results-of-the-bounded-final-exploration}}

The fully normalized original bundle has \textbf{1,144 fine AX
terminals}, compared with 1,270 originally and 1,162 after all six
direct double-negation deletions. It uses six DNE deletions plus one
additional negative-to-negative contraposition inside CP. It retains all
eight original logical rows, GEN, input reads, gate, cleanup, and the
final outer implication tag. Its exact data are in
\nolinkurl{boolean_first_variant.json}; the fine word alone is
\nolinkurl{boolean_first_variant.txt}.

The source shortening did not produce a new sub-245-state table during
this branch. The supplied full-normalization grammar seed needed 142
decision rules before further search, illustrating why primitive source
length is not a state-count result. The root report records the actual
finished machine candidates and must be used for their final state and
audit status.

A separate small family stays near original-layout DNE mask32, which
changes only redex ID5. \nolinkurl{boolean_local_variants.json} contains
the mask32 control and 31 distinct \textbf{single local Boolean edits},
for 32 targets total. The shortest has 1,234 fine AX terminals. These
are exact source targets for future grammar optimization; no new
transition-table count is assigned to them.

\hypertarget{logical-justification}{%
\subsection{Logical justification}\label{logical-justification}}

The common-core theorem in \nolinkurl{dne_calculus_transfer.md} supplies
MP, K, S, Peirce and ex falso for the actual numerical implication
operation and falsity code. Therefore every instance of a classical
implicational tautology is a common-core theorem. The six schemas used
by the additional generator are:

\begin{longtable}[]{@{}lll@{}}
\toprule\noalign{}
Rule & Left expression & Right expression \\
\midrule\noalign{}
\endhead
\bottomrule\noalign{}
\endlastfoot
Double negation & \texttt{¬¬A} & \texttt{A} \\
Negative contraposition & \texttt{¬A\ →\ ¬B} & \texttt{B\ →\ A} \\
Antecedent exchange & \texttt{A\ →\ (B\ →\ C)} &
\texttt{B\ →\ (A\ →\ C)} \\
Negative antecedent exchange & \texttt{¬A\ →\ B} & \texttt{¬B\ →\ A} \\
Negative consequent exchange & \texttt{A\ →\ ¬B} & \texttt{B\ →\ ¬A} \\
Unpack negative implication & \texttt{¬(A\ →\ B)\ →\ C} &
\texttt{A\ →\ (¬B\ →\ C)} \\
\end{longtable}

Here every negation means implication to the same numerical
\texttt{F=1}. All rows are equivalences, so both directional
implications are classical tautologies. The propositional completeness
theorem for arbitrary implication operations applies to their numerical
instances. The GEN-plus-rho2 congruence lemma then transports the
equivalence through any surrounding quantifiers, including a binder of
variable0. No binding position, parameter identity, source gate, or
schema domain changes.

Replace the old bundle axiom by any one of the checked new bundles. The
same two proof inductions give equality of the entire numerical theorem
sets. The explicit enumeration soundness invariant, derivation compiler,
and cofinal padding proof therefore carry over exactly as in the DNE
note. The new bundle still has an outer implication and can never be the
literal halt code1; this preserves the stack-selector proof's tag
premise.

\hypertarget{what-the-executable-checker-establishes}{%
\subsection{What the executable checker
establishes}\label{what-the-executable-checker-establishes}}

The standalone command

\begin{verbatim}
python boolean_bundle_variants.py --check
\end{verbatim}

reads the pinned source formula specification and reproduces the mask32
expression. It checks each of the six abstract Boolean equivalences on
all eight truth assignments to its three metavariables. For each
supplied local target, it then replays the stated rule at the stated
formula-position address, comparing the exact before and after subtrees.
It reconstructs the full emitted AX word, verifies all unchanged logical
rows and input events, checks the protected stack tail and outer
implication tag, and compares the exact fine word and hash. It also
independently replays all seven steps in the full 1,144-terminal target.

To regenerate the 32 local targets, run

\begin{verbatim}
python boolean_bundle_variants.py --generate
\end{verbatim}

The local target ledger records a zlib byte count only as a sorting
heuristic. That field is not a grammar-state estimate, a source
certificate, or a claim of machine size.

The checker uses \texttt{check\_dne\_variants.py},
\texttt{macro\_to\_fine.json}, \nolinkurl{boolean_first_variant.json},
and the small pinned \nolinkurl{snapshot/ZF/m254/tools/logic/spec.py}
fixture. It has no external Python package dependency. Its finite
truth-table calculation validates the rule schemas; the uniform
numerical proof relies on the common-core implication theorem stated and
proved in the mathematical note.

\hypertarget{grounding-the-fine-token-source-semantics}{%
\subsection{Grounding the fine-token source
semantics}\label{grounding-the-fine-token-source-semantics}}

The fine words refer to the routed baseline's optimized structured
source, whose internal increment order and control association can
differ from the original Wade macros. The equality of one complete
original AX string is insufficient to justify arbitrary new compositions
of its substrings.

The independent \nolinkurl{audit_macro_fine_bridge.py} resolves this
specific issue. For each of the 19 macros appearing in the emitted
words, it compares the original structured macro with the proposed fine
expansion at the level of finite control-flow graphs. Between successive
decrement guards it matches all nine accumulated register increments; it
checks both outcomes of each decrement and the exact return. It rechecks
all graph edges with the upstream path semantics. All 19 macros pass,
with 192 paired states and 346 edges. These are all-register equations,
not tests at sampled small values.

Consequently each symbolic RPN constructor step in the new word has the
original macro's arithmetic effect, under the original macro contract's
preconditions. Induction over the emitted word gives the exact changed
bundle and unchanged AX interface. This connects the Boolean/calculus
argument to the actual new structured source. The root's complete binary
coverage and exact blank frame then connect that source to its literal
transition table. These distinct compositions are all required; neither
the truth tables nor the macro bridge alone is a complete machine
consistency proof.

\hypertarget{attribution-and-status}{%
\subsection{Attribution and status}\label{attribution-and-status}}

The underlying compact calculus, set-theoretic calibration and original
constructor contracts are inherited and attributed as in
\nolinkurl{dne_calculus_transfer.md}. Prior incoming work already used
the same common-core DNE transport for a quantifier-row replacement.
This branch contributes fixed-bundle applications, exact altered source
words, local rewrite certificates, measured complete controls, and the
new source-interface composition.

No new Lean theorem build, unconditional nonhalting proof, minimality
result, or 244-state candidate is claimed in these supplementary files.
The full machine ledger belongs to the root report.
